%% LyX 2.4.4 created this file.  For more info, see https://www.lyx.org/.
%% Do not edit unless you really know what you are doing.
\documentclass[oneside,english]{book}
\usepackage[T1]{fontenc}
\usepackage[latin9]{inputenc}
\setcounter{secnumdepth}{3}
\setcounter{tocdepth}{3}
\usepackage{units}
\usepackage{amsmath}
\usepackage{amsthm}
\usepackage{amssymb}
\usepackage{geometry}
\geometry{verbose,tmargin=1in,bmargin=1in,lmargin=1.5in,rmargin=1.5in}

\makeatletter

%%%%%%%%%%%%%%%%%%%%%%%%%%%%%% LyX specific LaTeX commands.
%% Because html converters don't know tabularnewline
\providecommand{\tabularnewline}{\\}

%%%%%%%%%%%%%%%%%%%%%%%%%%%%%% Textclass specific LaTeX commands.
\theoremstyle{plain}
\newtheorem{thm}{\protect\theoremname}
\theoremstyle{definition}
\newtheorem{xca}{\protect\exercisename}
\theoremstyle{definition}
\newtheorem{example}{\protect\examplename}
\ifx\proof\undefined\
  \newenvironment{proof}[1][\proofname]{\par
    \normalfont\topsep6\p@\@plus6\p@\relax
    \trivlist
    \itemindent\parindent
    \item[\hskip\labelsep
          \scshape
      #1]\ignorespaces
  }{%
    \endtrivlist\@endpefalse
  }
  \providecommand{\proofname}{Proof}
\fi

%%%%%%%%%%%%%%%%%%%%%%%%%%%%%% User specified LaTeX commands.
\usepackage{hyperref}
\usepackage[usenames,dvipsnames]{xcolor} %adds additional color options
\hypersetup{colorlinks=true, urlcolor=Blue}%Blue

\usepackage{multicol}

\usepackage{tikz}
\usepackage{tikz-3dplot}
\usetikzlibrary{calc,arrows,shapes,backgrounds}

\usetikzlibrary{patterns}
\usetikzlibrary{plotmarks}
\usepackage{pgfplots}
\pgfplotsset{compat=newest}

\usepackage{calc}
\usepackage[nomessages]{fp}

\usepackage{datetime}
%\usepackage{subcaption}

%\usepackage{amsthm}
\usepackage[all]{xy}
% Added by lyx2lyx
\setlength{\parskip}{\smallskipamount}
\setlength{\parindent}{0pt}

\makeatother

\usepackage{babel}
\providecommand{\examplename}{Example}
\providecommand{\exercisename}{Exercise}
\providecommand{\theoremname}{Theorem}

\begin{document}

\chapter{Factor Rings}

\section{Properties of Factor Rings}

We saw in the previous chapter that $\mathbb{Z}/6\mathbb{Z}=\mathbb{Z}/\left\langle 6\right\rangle $
is isomorphic to $\mathbb{Z}_{6}$. It follows that $\mathbb{Z}/\left\langle 6\right\rangle $
is a ring. This is in general true for all factor rings, however we
have yet to prove this. Once we show that they are rings, we can then
explore their structure using isomorphisms and homomorphisms. First
we need to define the ring operations of factor rings which requires
the following theorem.

\emph{Recall definition \ref{def-factor-ring}:} If $I$ is an ideal
in the ring $R$, then the set of all cosets of $I$ is called the
\emph{factor ring} $R$ modulo $I$ and is denoted $R/I$.
\begin{itemize}
\item $R/I=\{r+I\colon r\in R\}$
\end{itemize}
$R/I$ is called a factor ring because (no surprise) -it is a ring
under addition and multiplication of modulo classes. Before we can
show that it is a ring, or why $I$ needs to be an ideal; we define
the operations and provide some examples. 
\begin{itemize}
\item The sum of $a+I$ and $b+I$ is $(a+b)+I$.

\begin{itemize}
\item []$a+I\oplus b+I=(a+b)+I$

\begin{itemize}
\item Don't forget that $x\in a+I\Rightarrow x=a+ri$ for some $i\in I$.
The addition sign can be confusing here because, $a+I$ denotes a
set -not an operation. Some texts write $a+I\oplus b+I$ as ``simply''$a+I+b+I$.
\end{itemize}
\end{itemize}
\item The product of $a+I$ and $b+I$ is $ab+I$

\begin{itemize}
\item []$a+I\otimes b+I=ab+I$ or $(a+I)(b+I)=ab+I$
\end{itemize}
\end{itemize}
The symbols $\oplus$ and $\otimes$ are used to distinguish the factor
ring operation from the coset notation. Though often used it is by
no means standard. So be careful when reading different texts! The
theorem below is just theorem \ref{thm:cong-prop} restated for ideals.
\begin{thm}
\label{thm:cong-for-mod-I}Let $I$ be an ideal in a ring $R$. If
$a+I=b+I$ and $c+I=d+I$ in $R/I$, then 

\[
(a+c)+I=(b+d)+I\mbox{ and }ac+I=bd+I
\]

\end{thm}
\begin{xca}
$\bigstar$ Prove the theorem above.
\end{xca}
\begin{example}
\label{ex-Z/<5> is a ring}$\left\langle 5\right\rangle $ is a principal
ideal in $\mathbb{Z}$. Addition and multiplication of cosets is just
the same as addition and multiplication of congruence classes $\mathbb{Z}$
modulo. 

$\mathbb{Z}/\left\langle 5\right\rangle =\left\{ 0+5\mathbb{Z},1+5\mathbb{Z},2+5\mathbb{Z},3+5\mathbb{Z},4+5\mathbb{Z}\right\} $.
We can add and multiply as follows: 
\[
\left(2+5\mathbb{Z}\right)\oplus\left(1+5\mathbb{Z}\right)=3+5\mathbb{Z}
\]
and 
\[
\left(3+5\mathbb{Z}\right)\otimes\left(2+5\mathbb{Z}\right)=6+5\mathbb{Z}=(5)+\left(1+5\mathbb{Z}\right)=\left(0+5\mathbb{Z}\right)\oplus\left(1+5\mathbb{Z}\right)=1+5\mathbb{Z}
\]

\begin{itemize}
\item []%
\begin{tabular}{rl}
$a+I\oplus b+I=(a+b)+I$ & closed under addition\tabularnewline
$a+I\oplus(b+I\oplus c+I)=(a+b+c)+I$ & additively associative\tabularnewline
$a+I\oplus0+I=a+I$ & additive identity\tabularnewline
$a+I\oplus-a+I=0+I$ & additive inverse\tabularnewline
\end{tabular}
\item []%
\begin{tabular}{rl}
$a+I\otimes b+I=ab+I$ & closed under multiplication\tabularnewline
$a+I\otimes(b+I\otimes c+I)=abc+I$ & multiplicatively associative\tabularnewline
$a+I\otimes(b+I\oplus c+I)=a+I\otimes((b+c)+I)$ & distribution\tabularnewline
$=\left(a(b+c)\right)+I=(ab+ac)+I$ & \tabularnewline
\end{tabular}So $R/I$ is a ring.
\end{itemize}
Thus $\mathbb{Z}/\left\langle 5\right\rangle =\mathbb{Z}/5\mathbb{Z}$
is a ring. To see that these are the distinct cosets, consider $a+5\mathbb{Z}$
for some integer $a>5$. $a=5q+r$ by the division algorithm so $a\equiv r\bmod5$
for some integer $0\leq r<5$. Thus $a$ is an element in one of the
above cosets, and it follows that $a+5\mathbb{Z}=r+5\mathbb{Z}$.
As the operations $\mathbb{Z}/\left\langle 5\right\rangle $ are just
modulo $5$ arithmetic, we can also see that $\mathbb{Z}/\left\langle 5\right\rangle =\mathbb{Z}_{5}$. 
\end{example}
\begin{xca}
Give two examples of elements in $\mbox{\ensuremath{\mathbb{Z}}}/4\mathbb{Z}=\mathbb{Z}/\left\langle 4\right\rangle $,
and add and multiply them.
\end{xca}

\begin{xca}
Find all the distinct cosets of $\mbox{\ensuremath{\mathbb{Z}}}/4\mathbb{Z}$.
\end{xca}
\begin{example}
$3\mathbb{Z}/9\mathbb{Z}=\left\{ 0+9\mathbb{Z},3+9\mathbb{Z},6+9\mathbb{Z}\right\} $.
For example $300=33\cdot9+3$ so $300\equiv3\bmod9$ and $300\in0+9\mathbb{Z}$. 

\label{ex-F=00005Bx=00005D/<p(x)> is a ring}If $F$ is a field and
$p(x)$ is a polynomial in $F[x]$, and $I$ is the principal ideal
$\left\langle p(x)\right\rangle $, then the cosets of $I$ are the
congruence classes modulo $p(x)$. Thus addition and multiplication
of cosets is done exactly as in chapter \ref{chap:Ideals}, and we
have that $F[x]/I$ is the congruence class ring $F[x]/\left\langle p(x)\right\rangle $.
\end{example}

\begin{example}
\label{ex-D/E is a ring}Let $D$ be the ring of rational numbers
(in reduced terms) with odd denominators, and let $E$ be the set
of elements in $D$ with even numerators. As we saw in exercise \ref{exc-D/R-factor-ring-of-even-denom},
$D/E$ has exactly two distinct ideals. $0+E$ and $x+E$ where $x$
is a rational number with an odd denominator and odd numerator. Since
$x-1=\frac{2n+1}{2m+1}-\frac{2m+1}{2m+1}=\frac{2(n-m)}{2m+1}\in I$.
So then $x+I=1+I$. So the two distinct cosets are $0+E$ and $1+E$.
By the theorem \ref{thm:cong-for-mod-I},\\
\begin{tabular}{c|cc}
$\oplus$ & $0+E$ & $1+E$\tabularnewline
\hline 
$0+E$ & $0+E$ & $1+E$\tabularnewline
$1+E$ & $1+E$ & $0+E$\tabularnewline
\end{tabular} and %
\begin{tabular}{c|cc}
$\otimes$ & $0+E$ & $1+E$\tabularnewline
\hline 
$0+E$ & $0+E$ & $0+E$\tabularnewline
$1+E$ & $0+E$ & $1+E$\tabularnewline
\end{tabular} so we can see that $D/E$ is also a ring.
\end{example}
\begin{thm}
\label{thm:R/I-is-a-ring}Let $I$ be an ideal in a ring $R$ then,

\begin{enumerate}
\item $R/I$ is a ring with addition and multiplication as defined above.
\item If $R$ is abelian then $R/I$ is also abelian.
\item If $R$ has an identity then $R/I$ has an identity.
\end{enumerate}
\end{thm}
\begin{proof}
$(1)$ First we need to show that $R/I$ is an abelian group under
addition. Recall from exercise \ref{exc-R/R=00003D0 and R/0=00003DR}
that $R/0=R$. $R$ is a ring so then $R/0$ is also a ring. We show
that $R/I$ is a subring of $R/0$. $R/I=\{r+I\colon r\in R\}\neq\emptyset$
since $I\neq\emptyset$ and $R\neq\emptyset$ as they are both rings.
Now let $a+I,b+I\in R/I$. $R$ is a ring; so $-b\in R$, $a-b\in R$,
and $ab\in R$. Thus $a+I\oplus-b+I=(a-b)+I\in R/I$, and $a+I\otimes b+I=ab+I\in R/I$
. So $R/I$ is subring (and hence a ring) by the subring test. 

Alternatively, we could show it is a ring directly as in example \ref{ex-Z/<5> is a ring}.

$(2)$ Suppose that $R$ is abelian so that $ab=ba$ for all $a,b\in R$.
This implies that $a+I\otimes b+I=ab+I=ba+I=b+I\otimes a+I$.

$(3)$ Suppose that $R$ has an identity, say $1_{R}$. This implies
that $a+I\otimes1_{R}+I=a1_{R}+I=a1_{R}+I=1_{R}+I\otimes a+I$. So
the identity of $R/I$ is $1_{R}+I$. 
\end{proof}
\begin{itemize}
\item Part 2 is a sufficient but necessary condition for $R/I$ to be commutative.
For example $R$ may not be non-abelian, but $R/R=0_{R}$ is certainly
abelian. 
\end{itemize}
\begin{thm}
\label{thm:R/I is a ring iif I is an ideal}Let $R$ be a ring. $R/I$
is a ring if and only if $I$ is an ideal of $R$. 
\end{thm}
\begin{proof}
$\left(\Leftarrow\right)$ From theorem \ref{thm:R/I-is-a-ring} $I$
is an ideal $\Rightarrow R/I$ is a ring.

$\left(\Rightarrow\right)$ Let $R/I$ be a ring and $I$ a subring
of $R$, but $I$ is not an ideal of $R$. Then there exists $a\in I$
such that $ra\notin I$ or $ar\notin I$ for some $r\in R$. Assume
that $ra\notin I$. Since $a\in I$, $a+I=0+I$ (cosets are equal
or disjoint). Now consider $r+I$. $(r+I)(a+I)=ra+I$, but $(r+I)(a+I)=(r+I)(0+I)=0+I=I$.
So then $ra+I=0+I$, a contradiction since $ra\notin I$. Similarly
we get a contradiction if $ar\notin I$. So multiplication is not
well defined and the set of cosets is not a ring.

Thus any ideal $I$ in $R$ gives us a factor ring by theorem \ref{thm:R/I-is-a-ring},
and any subring $A$ which is not an ideal gives us a non-ring $R/A$.
\end{proof}
\begin{xca}
Let $S=\left\{ a+bi\colon b\in2\mathbb{Z}\right\} $ ($b$ is even
and $i=\sqrt{-i}$). Prove that $S$ is a subring of $\mathbb{Z}[i]$
(see exercise \ref{ex-gaussian int factor ring }) but that $S$ is
not an ideal of $\mathbb{Z}[i]$. 
\end{xca}

\begin{xca}
Let $S=\left\{ a+bi\colon b\in2\mathbb{Z}\right\} $ ($b$ is even
and $i=\sqrt{-i}$). Prove that $\mathbb{Z}[i]/S$ is not a ring.
\end{xca}

\begin{xca}
Prove that $\mathbb{R}[x]/\left\langle x^{2}+1\right\rangle $ is
a ring. (We know that the ring of polynomials with integer coefficients,
$\mathbb{R}[x]$, is a ring). 
\end{xca}
\begin{example}
\label{ex-gaussian int factor ring }Let $\left\langle 4-i\right\rangle \in\mathbb{Z}[i]$.
$\mathbb{Z}[i]=\left\{ x+yi\colon x,y\in\mathbb{Z}\right\} $ is called
the \emph{Gaussian integers}, and\emph{. $\left\langle 4-i\right\rangle =\left\{ (4-i)\cdot z\colon z\in\mathbb{Z}[i]\right\} $.
}For any two Gaussian integers\emph{ }$z$\emph{ }and $z'$ we have
\emph{$(4-1i)z-(4-i)z'=(4-i)(z-z')\in\left\langle 4-i\right\rangle $},
and $z_{1}\left(z_{2}(4-i)\right)=\left(z_{2}\left(4-i\right)\right)z_{1}=z_{1}z_{2}\left(4-i\right)\in\left\langle 4-i\right\rangle $.
So $\left\langle 4-i\right\rangle $ is an ideal of $\mathbb{Z}[i]$. 

By theorem\ref{thm:R/I is a ring iif I is an ideal}, $\mathbb{Z}[i]/\left\langle 4-i\right\rangle $
is a ring with elements (cosets) of the form $x+yi+\left\langle 4-i\right\rangle $
for $x,y\in\mathbb{Z}$. When dealing with cosets in $\mathbb{Z}[x]/\left\langle 4-i\right\rangle $
we can treat $4-i$ as $0$ since $4-i+\left\langle 4-i\right\rangle =0+\left\langle 4-i\right\rangle $.
This allows us to simplify the set of \emph{distinct }cosets. For
example, since $4-i\equiv0$ we also know that $4\equiv i$ since
$4-i+i=4$. So then 
\[
6+3i+\left\langle 4-i\right\rangle =6+12+\left\langle 4-i\right\rangle =18+\left\langle 4-i\right\rangle 
\]

$4\equiv i\Rightarrow16\equiv-1$ and $17\equiv0$. We can now further
reduce the above coset to
\[
18+\left\langle 4-i\right\rangle =17+1+\left\langle 4-i\right\rangle =1+\left\langle 4-i\right\rangle 
\]
As it turns out, there are $17$ distinct cosets. 
\[
\mathbb{Z}[i]/\left\langle 4-i\right\rangle =\left\{ 0+\left\langle 4-i\right\rangle ,1+\left\langle 4-i\right\rangle ,\dots,17+\left\langle 4-i\right\rangle \right\} 
\]
To see that this is the case, note that any other coset will be some
integer multiple of $1+\left\langle 4-i\right\rangle $. Since $17\left(1+\left\langle 4-i\right\rangle \right)=0+\left\langle 4-i\right\rangle $,
$17\left(1+\left\langle 4-i\right\rangle \right)$ has order $1$
or $17$. If it is $1$, then $1+\left\langle 4-i\right\rangle =0+\left\langle 4-i\right\rangle $
and $1\in\left\langle 4-i\right\rangle $ so $1=\left(x+yi\right)\left(4-i\right)=4x-ix+4yi+y=4x+y+i(-x+4y)$
for integers $x$ and $y$. Since $1\in\mathbb{R}$ this implies that
$1=4x+y$ and $0=-x+4y$. Solving this system gives $y=\nicefrac{1}{17}\notin\mathbb{Z}$,
a contradiction. So the order of $\mathbb{Z}[x]/\left\langle 4-i\right\rangle $
is $17$. From the operations we can also see that this ring is essentially
the same as $\mathbb{Z}_{17}$, a field. 
\end{example}
\begin{xca}
Give two examples of elements in $\mbox{\ensuremath{\mathbb{Z}}}[i]/\left\langle 2-i\right\rangle $,
and add and multiply them.
\end{xca}

\begin{xca}
Find all the distinct cosets of $\mbox{\ensuremath{\mathbb{Z}}}[i]/\left\langle 2-i\right\rangle $.
\end{xca}


\section{Prime and Maximal Ideals }
\begin{xca}
Let $R$ be a commutative ring with unity and let $a_{1},a_{2},\dots,a_{n}\in R$.
Prove that 
\[
I=\left\langle a_{1},a_{2},\dots,a_{n}\right\rangle =\left\{ r_{1}a_{1},r_{2}a_{2},\dots,r_{n}a_{n}\colon r_{i}\in R\right\} 
\]
 is an ideal in $R$. $I$ is called the \emph{ideal generated by
}$a_{1},a_{2},\dots,a_{n}$. 
\end{xca}

\begin{xca}
If $n$ is not prime prove that $\left\langle n\right\rangle $ is
not a prime ideal in $\mathbb{Z}$
\end{xca}

\begin{xca}
Prove that the intersection of two prime ideals need not be prime.\\
\end{xca}

\begin{xca}
Prove that a nonzero integer $p$ is prime if and only if the ideal
$\left\langle p\right\rangle $ is maximal in $\mathbb{Z}$.
\end{xca}

\begin{xca}
Let $F$ be a field and $p(x)\in F[x]$. Prove that $p(x)$ is irreducible
if and only if the ideal $\left\langle p(x)\right\rangle $ is maximal
in $F[x]$.
\end{xca}
\begin{itemize}
\item All ideals of $\mathbb{Z}_{n}$ are principal. 
\item The ideals of $\mathbb{Z}_{n}$ are the divisors of $n$. To see this
let $\left\langle m\right\rangle $ be an ideal in $\mathbb{Z}_{n}$.
$0_{R}\in\left\langle m\right\rangle $ so then there is a $r\in R$
such that $rm=n\equiv0\bmod n$. 
\item The distinct divisors of $n$ give distinct ideals of $\mathbb{Z}_{n}$.
To see this let $0<m<q$ be two distinct divisors of $n$. If $m|q$
then $\left\langle q\right\rangle \subset\left\langle m\right\rangle $.
Otherwise $m\notin\left\langle q\right\rangle $.
\end{itemize}
\begin{example}
Consider $\mathbb{Z}_{36}$. The distinct ideals are $\left\langle 0\right\rangle $,
$\left\langle 12\right\rangle $, $\left\langle 18\right\rangle $,
$\left\langle 4\right\rangle $, $\left\langle 6\right\rangle $,
\textbf{$\left\langle 9\right\rangle $}, $\left\langle 2\right\rangle $,
$\left\langle 3\right\rangle $, and $\left\langle 1\right\rangle =\mathbb{Z}_{36}$.

\begin{itemize}
\item $\left\langle 0\right\rangle \subset\left\langle 18\right\rangle \subset\left\langle 9\right\rangle \subset\left\langle 3\right\rangle \subset\left\langle 1\right\rangle $
\item $\left\langle 0\right\rangle \subset\left\langle 18\right\rangle \subset\left\langle 6\right\rangle \subset\left\langle 3\right\rangle \subset\left\langle 1\right\rangle $
\item $\left\langle 0\right\rangle \subset\left\langle 18\right\rangle \subset\left\langle 6\right\rangle \subset\left\langle 2\right\rangle \subset\left\langle 1\right\rangle $
\end{itemize}
Continuing like this we can construct a lattice showing all the ideals
and clearly see that $\left\langle 2\right\rangle $ and $\left\langle 3\right\rangle $
are the maximal ideals. 
\end{example}
\begin{xca}
Find all maximal ideals in $\mathbb{Z}_{10}$ and $\mathbb{Z}_{12}$.
\end{xca}
\begin{thm}
\label{thm:I is prime <=00003D> R/I is an integral domain}Let $R$
be a commutative ring with identity. $R/I$ is an integral domain
if and only if $I$ is a prime ideal.

\begin{proof}
$(\Rightarrow)$ Suppose that $R/I$ is an integral domain and $ab\in I$.
Then $(a+I)(b+I)=ab+I=I=0+I$. So then $ab$ equals the zero element
of $R$. $R$ is an integral domain so then $a=0$ or $b=0$. Thus
$a\in I$ or $b\in I$ implying that $I$ is prime.

$(\Leftarrow)$ Suppose that $I$ is a prime ideal. $R$ is commutative
with identity, so we need only show that $R$ has no zero divisors
(nonzero elements $a$ such that $a\cdot b=0_{R/I}$ for any nonzero
$b$). Suppose that $(a+I)(b+I)=ab+I=0+I=I$. So then $ab\in I$ and
thus $a\in I$ or $b\in I$ since $I$ is prime. So then $a+I$ OR
$b+I$ is equal to $0+I$, and $R/I$ is an integral domain. 
\end{proof}
\end{thm}

\begin{thm}
\label{thm:I is max<=00003D>R/I  is field}Let $R$ be a commutative
ring with identity. $R/I$ is a field if and only if $I$ is a maximal
ideal.

\begin{proof}
$(\Rightarrow)$ Let $R/I$ be a field, and let $M$ be an ideal of
$R$ such that $I\subset M$ (properly contains). We want to show
that $M=R$. So there is a $m\in M$ such that $m\notin I$. So the
$m+I\neq0+I$ and because $R/I$ is a field there exists an $x\in R$
such that $(m+I)(x+I)=mx+I=1+I$. 

Since $m\in M$ we have that $mx\in M$ ($M$ is an ideal), and so
$mx=1\in M$. By a previous exercise this implies that $M=R$. So
$I$ is maximal.

$(\Leftarrow)$ Let $I$ be a maximal ideal. Let $m\in R$ but $m\notin I$.
We want to show that $m+I$ must have a multiplicative inverse. Consider
$M=\left\{ mr+i\colon r\in R\mbox{ and }i\in I\right\} $. It is easy
to see that $M$ is an ideal in $R$, and $I\subset R$. $I$ is maximal
so then $M=R$. Thus $1\in M$, say $1=mx+i'$ where $i'\in I$. Then
\[
1+I=(mx+i')+I=(mx+I)\oplus(i'+I)=(mx+I)\oplus(0+I)=mx+I
\]
So then $mx+I=\left(m+I\right)\otimes(x+I)$. Thus $m+I$ has an inverse
and $R/I$ is a field. 
\end{proof}
\end{thm}
\begin{itemize}
\item A field is an integral domain, so theorem \ref{thm:I is max<=00003D>R/I  is field}
tells us that a maximal ideal is always prime. Of course a prime ideal
need not be maximal as the example below shows.
\end{itemize}
\begin{example}
$\left\langle x\right\rangle $ is a prime ideal in $\mathbb{Z}[x]$.
$\left\langle x\right\rangle =\left\{ f(x)\in\mathbb{Z}[x]\colon f(0)=0\right\} $,
zero constants. Thus if $g(x)h(x)\in\left\langle x\right\rangle $
then $g(0)h(0)=0$ and either $g(0)=0$ or $h(0)=0$. 

But $\left\langle x\right\rangle $ is not maximal since $\left\langle x\right\rangle \subset\left\langle x,2\right\rangle \subset\mathbb{Z}[x]$.
It follows that $\mathbb{Z}[x]/\left\langle x\right\rangle $ is an
integral domain but not a field. 
\end{example}

\section{Factor Rings and Homomorphisms}

Factor rings arose as a generalization of congruence class arithmetic
in the integers. Now we see that the concept of homomorphism and isomorphism
that we studied in chapter \ref{chap:Homomorphisms-and-Isomorphisms}
is closely related to factor rings and ideals. 
\begin{itemize}
\item Recall definition \ref{def: A-(ring)-homomorphism,}: a \emph{(ring)
homomorphism} is an order preserving function between rings, i.e.,
a homomorphism between rings $R$ and $S$ is a function, say $f$
written $f\colon R\rightarrow S$, such that $f(ab)=f(a)f(b)$ and
$f(a+b)=f(a)+f(b)$.
\end{itemize}
\begin{thm}
\label{thm:kernel is an ideal}Let $f\colon R\rightarrow S$ be a
ring homomorphism and let $\ker f=\left\{ r\in R\colon f(r)=0_{S}\right\} $.
Then $\ker f$ is an ideal in $R$. 

\begin{proof}
$f$ is a homomorphism so $f(0_{R})=0_{S}$ by theorem \ref{thm:Properties-of-Ring},
and we have that $\ker f\neq\emptyset$. Using the ideal test (theorem
\ref{thm:Ideal-test.-A}), let $k,k'\in\ker f$. $R$ is homomorphism
so $f(k-k')=f(k)-f(k')=0_{S}-0_{S}=0_{S}\Rightarrow r-r'\in\ker f$.
Now for any$r\in R$ the homomorphic properties of $f$ implies that
$f(rk)=f(r)f(k)=f(r)0_{S}=0_{S}\Rightarrow rk\in\ker f$. Similarly
we can show that $kr$ is in $\ker f$. So $\ker f$ is an ideal in
$R$.
\end{proof}
\end{thm}
\begin{itemize}
\item The ideal $\ker f$ from theorem \ref{thm:kernel is an ideal} is
called the \emph{kernel of }$f$. It is simply the set of elements
in $R$ that is mapped to $0$. 
\end{itemize}
\begin{example}
Recall that the floor function (see example \ref{ex-floor-funct})
$f\colon\mathbb{R\rightarrow\mathbb{Z}}$ is a homomorphism. $\ker f=\left\{ x\in\mathbb{Z}\colon0\leq x<1\right\} $. 
\end{example}
\begin{xca}
\label{exc: constant term of poly map}Show that $\sigma\colon\mathbb{Z}[x]\rightarrow\mathbb{Z}$
defined by $\sigma(a_{n}x^{n}+a_{n-1}x^{n-1}\dots+a_{1}x+a_{0})=a_{0}$
is a homomorphism and find $\ker f$.
\end{xca}

\begin{xca}
\label{exc: ker or Z->Zn}In exercise \ref{exc-z->zn is a homo} we
showed that $f\colon\mathbb{Z}\rightarrow\mathbb{Z}_{n}$ given by
$k\rightarrow k\mod n$ is a ring homomorphism. Find $\ker f$.
\end{xca}

\begin{xca}
$\bigstar$ Show that the $f\colon\mathbb{Z}_{12}\rightarrow\mathbb{Z}_{4}$
defined by $f(x)=x\bmod4$ is a onto homomorphism, and find $\ker f$.
\end{xca}

\begin{xca}
\label{exc: ker exc with matrix to a-b}Let $R=\left\{ \begin{bmatrix}\begin{array}{rr}
a & b\\
c & d
\end{array}\end{bmatrix}\colon a,b\in\mathbb{Z}\right\} $, and let $\phi\colon R\rightarrow\mathbb{Z}$ be defined by $\phi\left(\begin{bmatrix}\begin{array}{rr}
a & b\\
c & d
\end{array}\end{bmatrix}\right)=a-b$. Show that $\phi$ is a homomorphism.
\end{xca}

\begin{xca}
Let $\phi$ be the function defined above in exercise \ref{exc: ker exc with matrix to a-b}.
Find $\ker\phi$. 
\end{xca}

If $\ker f=R$, meaning all elements are mapped to $0$, then $f$
is the zero function, $f=0$. If $\ker f$ has more than one element
then $f(a)=f(b)=0$ for some $a\neq b$ which is the definition of
not being one-to-one (injective). So the kernel gives us a sense of
how close a function is to being one-to-one which leads us to the
following theorem.
\begin{thm}
\label{thm:kerf=00003D0 iff injective}Let $f\colon R\rightarrow S$
be a ring homomorphism with kernel $\ker f$. Then $\ker f=\left\{ 0_{R}\right\} $
if and only if $f$ is one-to-one. 

\begin{proof}
$(\Leftarrow)$ Suppose that $f$ is one-to-one so that $f(a)=f(b)\iff a=b$
for any $a,b\in R$. Since $f$ is a homomorphism $f(0_{R})=0_{S}$
by theorem \ref{thm:Properties-of-Ring}. Because $f$ is one-to-one,
$0_{R},a\in\ker f\iff f(0_{R})=f(a)=0_{S}\iff0_{R}=a$. Thus $\ker f=\left\{ 0_{R}\right\} $.

$(\Rightarrow)$ Suppose that $\ker f=\left\{ 0_{R}\right\} $ and
$f(a)=f(b)$. We need to show that $a=b$. Because $f$ is a homomorphism,
$f(a-b)=f(a)-f(b)=0_{S}\Rightarrow a-b\in\ker f$. By assumption $a-b=0_{R}\Rightarrow a=b$. 
\end{proof}
\end{thm}
\begin{xca}
Use theorem \ref{thm:kerf=00003D0 iff injective} to determine whether
the functions from exercise \ref{exc: ker or Z->Zn} and \ref{exc: ker exc with matrix to a-b}
are one-to-one.
\end{xca}
Theorem \ref{thm:kernel is an ideal} tells us that every kernel is
an ideal. Conversely, every ideal is a kernel of some homomorphism.
This is where we connect factor rings and homomorphisms. Recall that
if $I$ is an ideal in $R$ then $R/I$ is also a ring. So we can
define a function from the ring $R$ to the ring $R/I$.
\begin{thm}
\label{thm:natural homo}Let $I$ be an ideal in the ring $R$. Then
the function $\pi\colon R\rightarrow R/I$ defined by $\pi(r)=r+I$
is an onto homomorphism and $\ker\pi=I$. 

\begin{proof}
It is easy to see that $\pi$ is onto since for any coset $r+I\in R/I$,
$\pi(r)=r+I$ by definition. The homomorphism part follows by the
definition of multiplication and addition in cosets. For any $r,r'\in R$
\[
f(r+r')=(r+r')+I=\left(r+I\right)\oplus\left(r'+I\right)
\]
 and 
\[
f(r\cdot r')=\left(r\cdot r'\right)+I=\left(r+I\right)\otimes\left(r'+I\right).
\]
The $\ker\pi$ will b e all the elements in $r\in R$ such that $\pi(r)=r+I=0+I=I$,
and $r+I=I$ if and only if $r\in I$. Thus $\ker\pi=I$.
\end{proof}
\end{thm}
\begin{itemize}
\item The function $\pi$ from theorem \ref{thm:natural homo} is called
the \emph{natural homomorphism}. 
\end{itemize}
\begin{xca}
Let $R$ and $S$ be rings. Show that $\pi\colon R\times S\rightarrow R$
defined by $\pi(r,s)=r$ is a onto homomorphism with $\ker f$ isomorphic
to $S$. 
\end{xca}
The natural homomorphism is a special case of when $f\colon R\rightarrow S$
is an onto homomorphism. When this happens we say that $f$ is \emph{homomorphic
image of $R$}. In the case that $f$ is also one-to-one it is an
isomorphism, and we know that $R$ and $S$ have identical structures.
Meaning that the algebraic properties in one ring will be identical
to those in the other ring. The next theorem relates the kernel of
a homomorphic image and the rings. 
\begin{thm}
\textbf{\emph{\label{thm:First-Isomorphism-Theorem.}First Isomorphism
Theorem. }}Let $f\colon R\rightarrow S$ be a onto homomorphism of
rings. Then the factor ring $R/\ker f$ is isomorphic to $S$.

\begin{proof}
Let $\phi\colon R/\ker f\rightarrow S$. We want to show that $\phi$
is a isomorphism. To define $\phi$ we need to associate with each
coset $r+\ker f\in R/\ker f$ an element of $S$. Let $\phi(r+\ker f)=f(r)$. 

For $\phi$ to be well defined (i.e., a function) we need to show
that if $r+\ker f=r'+\ker f$ then $f(r)=f(r')$.

\begin{itemize}
\item The problem is that a coset may possibly be represented by many different
elements in $R$
\end{itemize}
If $r+\ker f=r'+\ker f$ then $r-r'\in\ker f$ by since $\ker f$
is an ideal, and since $f$ is a homomorphism $f(r)-f(r')=f(r-r')=0_{S}$.
This $f(r)-f(r')=0_{S}$ implies that $f(r)=f(r')$ as needed. Thus
$\phi\colon R/\ker f\rightarrow S$ as defined by $\phi(r+\ker f)=f(r)$
is a function.

Now we need to show that $\phi$ is a isomorphism. If $s\in S$ then
$s=f(r)$ for some $r\in R$ because $f$ is onto. Then $s=f(r)=\phi(r+\ker f)$,
and $\phi$ is also onto. If $\phi(r+\ker f)=\phi(r'+\ker f)$ then
$f(r)=f(r')\Rightarrow f(r)-f(r')=0_{S}\Rightarrow f(r-r')=0_{S}$.
Hence $r-r'\in\ker f$ which implies that $r+\ker f=r'+\ker f$. Therefore
$\phi$ is also one-to-one. 

Using the definition of addition and multiplication of cosets along
with the homomorphic properties of $f$ we show that $\phi$ is a
homomorphism:
\begin{eqnarray*}
\phi\left[(r+\ker f)\oplus(r'+\ker f)\right] & = & \phi\left[(r+r')+\ker f\right]=f(r+r')=f(r)+f(r')\\
 & = & \phi(r+\ker f)\oplus\phi(r'+\ker f)
\end{eqnarray*}
and 
\begin{eqnarray*}
\phi\left[(r+\ker f)\otimes(r'+\ker f)\right] & = & \phi\left[r\cdot r'+\ker f\right]=f(r\cdot r')=f(r)f(r')\\
 & = & \phi(r+\ker f)\otimes\phi(r'+\ker f)
\end{eqnarray*}
Thus $\phi$ is an isomorphism.
\end{proof}
\end{thm}
\begin{figure}[h]
\begin{equation*}   
	\xymatrix@C+1em@R+1em{    
		R \ar[r]^-{f} \ar[d]_-{\pi} & S \\    
		R/\ker f \ar@{-->}[ur]_-{\tilde{f}}   
	}  
\end{equation*}

\caption{\label{fig:iso thm graph}A graph representing theorem \ref{thm:First-Isomorphism-Theorem.}.}

\end{figure}

This theorem states that every homomorphic image of a ring $R$ is
an isomorphism to $R/\ker f$ for some ideal $\ker f$. If you know
all the factor rings of $R$ then you also know all possible homomorphic
images of $R$. The ideal $\ker f$ measures how much information
is lost in the mapping from $R$ to $R/\ker f$. When $\ker f=\left\{ 0_{R}\right\} $,
no information is lost -it is an isomorphism. 
\begin{itemize}
\item Theorem \ref{thm:First-Isomorphism-Theorem.} implies that all homomorphic
images of $R$ can be determined using the factor rings of $R$.
\end{itemize}
The proof of theorem \ref{thm:First-Isomorphism-Theorem.} shows that
$\tilde{f}\circ\pi=f$ , i.e., the graph in figure \ref{fig:iso thm graph}
is commutative. 
\begin{example}
Recall that the principal ideal $\left\langle x\right\rangle $ in
the ring $\mathbb{Z}[x]$ consists of all polynomials with constant
term $0$ (polynomial multiples of $x$). To understand $\mathbb{Z}[x]/\left\langle x\right\rangle $
we can look at the function $\sigma\colon\mathbb{Z}[x]\rightarrow\mathbb{Z}$
defined as in example \ref{exc: constant term of poly map} (it gives
the constant term of a polynomial). $\sigma$ is onto as for any integer
$a\in\mathbb{Z}$, $a$ is also a polynomial with integer coefficients
so then $\sigma(a)=a$. In exercise \ref{exc: constant term of poly map}
it was shown that $\sigma$ is also a homomorphism. The kernel consists
of all the elements in $\mathbb{Z}[x]$ mapped to $0$. This is all
polynomials without a constant term; which is just $\left\langle x\right\rangle $.
So then $\ker\sigma=\left\langle x\right\rangle $. Thus by the first
isomorphism theorem, $\mathbb{Z}[x]/\ker\sigma\cong\mathbb{Z}$\footnote{Recall that ``$\cong$'' notates ``isomorphic to''.}.
\end{example}

\begin{example}
Let $S$ be a homomorphic image of $\mathbb{Z}$. We want to classify
all possible such images. If $S$ is a homomorphic image of $\mathbb{Z}$
then there is a homomorphic onto function $f\colon\mathbb{Z}\rightarrow S$.
If $f$ is a homomorphism then the structure of $S$ is exactly like
that of the integers. 

Now suppose that $S$ is not isomorphic to $\mathbb{Z}$. Then the
kernel of $f$ is a nonzero ideal, i.e., $\ker f\neq\left\langle 0\right\rangle $.
Because $\ker f$ is an ideal in $\mathbb{Z}$ it must be a principal
ideal (a previous example). So then $\ker f=\left\langle n\right\rangle $
for some integer $n$. So then by theorem \ref{thm:First-Isomorphism-Theorem.}
$S\cong\mathbb{Z}/\left\langle n\right\rangle $ which we know is
the same as $\mathbb{Z}_{n}$. So any homomorphic image of the integers
is either $\mathbb{Z}$ or $\mathbb{Z}_{n}$ for some $n\in\mathbb{Z}$.
\end{example}
\begin{xca}
Use theorem \ref{thm:First-Isomorphism-Theorem.} to show that $R/\left\langle 0_{R}\right\rangle \cong R$.
\end{xca}

\begin{xca}
Use theorem \ref{thm:First-Isomorphism-Theorem.} to show that $\mathbb{Z}_{20}/\left\langle 5\right\rangle \cong\mathbb{Z}_{5}$.
\end{xca}

\begin{xca}
Show that $R/\ker f\cong\mathbb{Z}$ where $R$ and $f$ are defined
in exercise \ref{exc: ker exc with matrix to a-b}.
\end{xca}

\begin{xca}
Is the kernel $\ker f$ from \ref{exc: ker exc with matrix to a-b}
a prime ideal? Is it a maximal ideal?
\end{xca}

\end{document}
